\begin{abstract}

% MANUSCRIPT-NODE: ABSTRACT-P01
Negative feedback is a central resource in policy optimization: it suppresses known bad behavior and can move a policy beyond the solution obtained from positive examples alone. Yet the same feedback can become destructive when historical negative actions remain in offline data, replay, or stale trajectories after the learner has moved away from them. We identify the missing mechanism by separating sample badness from learner-relative distance or rarity, showing that policy remoteness independently amplifies negative influence. We then develop Repulsive Dynamics, which characterizes a transition from the Positive-only target, through stable extrapolation, to a feasibility boundary and the loss of a finite equilibrium. Based on this analysis, Distributionally Robust Policy Optimization (DRPO) reweights the aggregate negative term with a policy-relative exponential envelope, retaining useful local repulsion while suppressing its far-field tail. Controlled continuous and categorical environments isolate the source and causal transmission of the effect; registered Hopper/D4RL and Countdown experiments are reserved for external validation under learned critics and shared-parameter sequence policies, with formal results remaining TBD until terminal audit. The resulting framework tests a general principle: policy optimization should control far-field negative influence rather than remove negative feedback as a whole.
% END-MANUSCRIPT-NODE: ABSTRACT-P01
\end{abstract}
